% ==============================================================================
% Capítulo 3: Soluções relacionadas
% ==============================================================================
\chapter{Soluções relacionadas}
\label{cap:solucoes}

\begin{instrucao}
\orientacao{O que vai em Soluções Relacionadas}
Em um trabalho tecnológico, este capítulo corresponde aos trabalhos
relacionados de um trabalho científico. As competências do TCC pedem que o
projeto seja relacionado às soluções de mercado equivalentes, demonstrando sua
relevância. Analise:
\begin{itemize}
  \item produtos e serviços que resolvem o mesmo problema ou parte dele;
  \item sistemas institucionais já existentes;
  \item trabalhos acadêmicos (artigos e TCCs) que desenvolveram sistemas
        semelhantes.
\end{itemize}
Defina critérios de comparação ligados aos requisitos do seu sistema e
aplique-os a todas as soluções. Termine com um quadro comparativo e um
parágrafo que mostre o que o seu sistema oferece de diferente. Se alguma
solução já resolve bem o problema, diga isso e explique por que ela não serve
ao seu contexto (custo, integração, privacidade dos dados).

Ao descrever produtos comerciais, informe a data de consulta, porque
funcionalidades e preços mudam. Cite a página do produto como referência
(\verb|@misc| com \verb|url| e \verb|urldate|).
\end{instrucao}

As soluções usadas hoje para organizar monitorias se dividem em ferramentas
genéricas de agendamento, planilhas e formulários compartilhados e o sistema
acadêmico institucional. Esta análise considerou os critérios que emergiram
das entrevistas com os usuários, descritas no \cref{cap:requisitos}.

% ------------------------------------------------------------------------------
\section{Ferramentas genéricas de agendamento}

Serviços de agendamento \emph{on-line} permitem que uma pessoa publique seus
horários livres e que outras reservem esses horários por um endereço web.
Eles evitam reservas simultâneas e enviam lembretes, mas são pensados para
atendimentos individuais de profissionais. Não associam a reserva a uma
disciplina, não registram a presença do estudante nem o assunto atendido e,
nas versões gratuitas, costumam limitar o número de agendas por conta. Além
disso, exigem que os estudantes informem dados pessoais a um serviço externo
à universidade.

% ------------------------------------------------------------------------------
\section{Planilhas e formulários compartilhados}

Planilhas e formulários são a solução mais usada pelos monitores entrevistados.
São gratuitos e conhecidos, mas não impedem que dois estudantes reservem o
mesmo horário, não enviam avisos e dependem do monitor para consolidar os
registros de frequência ao fim do mês. Os professores relataram que os
registros chegam incompletos e em formatos diferentes de um monitor para outro.

% ------------------------------------------------------------------------------
\section{Sistema acadêmico institucional}

O sistema acadêmico da universidade registra os monitores vinculados a cada
disciplina e controla o pagamento de bolsas, mas não oferece funcionalidades
de agendamento ou de registro de atendimentos. Integrar o Monitora a esse
sistema seria desejável, porém exigiria acesso a interfaces de programação
que não estão disponíveis para projetos de estudantes.

% ------------------------------------------------------------------------------
\section{Comparação}

O \cref{qua:comparacao} compara as três categorias com o sistema proposto. O
Monitora é a única opção que reúne agendamento sem conflitos, registro de
frequência e de assuntos atendidos e relatórios por disciplina, mantendo os
dados dos estudantes em infraestrutura da universidade.

\begin{quadro}[htbp]
  \centering
  \caption{Comparação entre as soluções existentes e o Monitora}
  \label{qua:comparacao}
  \small
  \begin{tabularx}{\textwidth}{|X|c|c|c|c|}
    \hline
    \textbf{Critério} & \textbf{Agendadores} & \textbf{Planilhas} &
      \textbf{\begin{tabular}[c]{@{}c@{}}Sistema\\acadêmico\end{tabular}} &
      \textbf{Monitora} \\
    \hline
    Impede reservas simultâneas & Sim & Não & Não & Sim \\
    \hline
    Associa atendimento à disciplina & Não & Parcial & Sim & Sim \\
    \hline
    Registra frequência e assunto & Não & Parcial & Não & Sim \\
    \hline
    Envia lembretes & Sim & Não & Não & Sim \\
    \hline
    Gera relatórios por disciplina & Não & Não & Parcial & Sim \\
    \hline
    Dados em infraestrutura da Ufes & Não & Não & Sim & Sim \\
    \hline
  \end{tabularx}
  \fonte{Elaborado pelo próprio autor com base nas entrevistas e em consultas
    realizadas em outubro de 2026.}
\end{quadro}

\begin{instrucao}
\orientacao{Quadro ou tabela?}
A NBR 14724 distingue os dois. O \textbf{quadro} organiza informações
textuais em linhas e colunas e tem bordas fechadas, com linhas verticais,
como o \cref{qua:comparacao}. A \textbf{tabela} apresenta dados, em geral
numéricos, e segue as normas do IBGE, com bordas laterais abertas e sem linhas
verticais (veja a \cref{tab:cobertura} no \cref{cap:testes}). Quadros e
tabelas têm listas e numeração próprias. A identificação fica acima e a fonte
abaixo, como nas figuras.
\begin{exemplocodigo}
\begin{quadro}[htbp]
  \centering
  \caption{Comparação entre as soluções}
  \label{qua:comparacao}
  \small
  \begin{tabular}{|l|c|c|}
    \hline
    \textbf{Critério} & \textbf{Planilhas} & \textbf{Monitora} \\
    \hline
    Impede reservas simultâneas & Não & Sim \\
    \hline
  \end{tabular}
  \fonteautor
\end{quadro}
\end{exemplocodigo}
Para montar tabelas e quadros mais elaborados, o site
\href{https://www.tablesgenerator.com/}{Tables Generator}
(\url{https://www.tablesgenerator.com/}) gera o código LaTeX a partir de uma
planilha. Depois de colar o código, ajuste-o ao padrão do modelo (ambiente
\verb|table| ou \verb|quadro|, \verb|\caption| acima e \verb|\fonte| abaixo).
\end{instrucao}
